% Shared preamble for all spec documents.
% Include from each spec: % Shared preamble for all spec documents.
% Include from each spec: % Shared preamble for all spec documents.
% Include from each spec: \input{../preamble.tex} (adjust depth).

\documentclass[11pt,a4paper]{article}

\usepackage[utf8]{inputenc}
\usepackage[T1]{fontenc}
\usepackage{lmodern}
\usepackage[a4paper,margin=2.3cm]{geometry}
\usepackage{parskip}
\usepackage{microtype}
\usepackage[dvipsnames]{xcolor}
\usepackage{enumitem}
\usepackage{listings}
\usepackage{tcolorbox}
\usepackage{titlesec}
\usepackage{fancyhdr}
\usepackage{array}
\usepackage{longtable}
\usepackage{booktabs}
\usepackage{amsmath}
\usepackage{amssymb}
\usepackage{hyperref}
\usepackage{cleveref}

\tcbuselibrary{skins,breakable}

\definecolor{specBlue}{RGB}{30,64,132}
\definecolor{specGray}{RGB}{90,90,90}
\definecolor{specAccent}{RGB}{198,40,40}
\definecolor{specGreen}{RGB}{30,110,60}

\hypersetup{
  colorlinks=true,
  linkcolor=specBlue,
  urlcolor=specBlue,
  citecolor=specBlue,
  pdfborder={0 0 0}
}

\titleformat{\section}{\Large\bfseries\color{specBlue}}{\thesection}{0.75em}{}
\titleformat{\subsection}{\large\bfseries\color{specBlue!80}}{\thesubsection}{0.6em}{}

\makeatletter
\newcommand{\specID}[1]{\def\@specID{#1}}
\newcommand{\specTitle}[1]{\def\@specTitle{#1}}
\newcommand{\specStatus}[1]{\def\@specStatus{#1}}
\newcommand{\specVersion}[1]{\def\@specVersion{#1}}
\newcommand{\specOwner}[1]{\def\@specOwner{#1}}
\newcommand{\specTraces}[1]{\def\@specTraces{#1}}

\specID{SPEC-XXX}
\specTitle{Untitled Spec}
\specStatus{DRAFT}
\specVersion{0.1}
\specOwner{Jeremie Nunez}
\specTraces{}

\newcommand{\makespecheader}{%
  \begin{center}
    {\LARGE\bfseries\color{specBlue}\@specTitle}\\[0.4em]
    {\small\color{specGray}\texttt{\@specID} \quad v\@specVersion \quad \textsc{\@specStatus} \quad \@specOwner}
  \end{center}
  \vspace{0.4em}
  \hrule
  \vspace{0.8em}
}

\pagestyle{fancy}
\fancyhf{}
\fancyhead[L]{\small\color{specGray}\@specID}
\fancyhead[R]{\small\color{specGray}\@specTitle}
\fancyfoot[C]{\small\color{specGray}\thepage}
\renewcommand{\headrulewidth}{0pt}
\makeatother

\newtcolorbox{invariant}{
  colback=specBlue!5, colframe=specBlue!75,
  title=Invariant, fonttitle=\bfseries, breakable,
  boxrule=0.5pt, arc=2pt
}

\newtcolorbox{scenario}[1]{
  colback=white, colframe=specBlue!50,
  title={Scenario: #1}, fonttitle=\bfseries, breakable,
  boxrule=0.5pt, arc=2pt
}

\newtcolorbox{ac}[1]{
  colback=specAccent!5, colframe=specAccent!60,
  title={Acceptance Criterion #1}, fonttitle=\bfseries, breakable,
  boxrule=0.5pt, arc=2pt
}

\newtcolorbox{nongoal}{
  colback=specGray!8, colframe=specGray!60,
  title=Non-Goal, fonttitle=\bfseries, breakable,
  boxrule=0.5pt, arc=2pt
}

\lstdefinestyle{codestyle}{
  basicstyle=\ttfamily\small,
  breaklines=true,
  showstringspaces=false,
  frame=single,
  framesep=4pt,
  rulecolor=\color{specBlue!30},
  backgroundcolor=\color{specBlue!2},
  xleftmargin=0.5em,
  xrightmargin=0.5em,
  keywordstyle=\color{specBlue}\bfseries,
  commentstyle=\color{specGray}\itshape,
  stringstyle=\color{specGreen}
}
\lstset{
  inputencoding=utf8,
  extendedchars=true,
  literate=
    {—}{{--}}1
    {–}{{-}}1
    {…}{{...}}1
    {’}{{'}}1
    {‘}{{'}}1
    {“}{{``}}1
    {”}{{''}}1
    {é}{{\'e}}1
    {è}{{\`e}}1
    {ê}{{\^e}}1
    {à}{{\`a}}1
    {â}{{\^a}}1
    {ç}{{\c{c}}}1
    {ô}{{\^o}}1
    {→}{{$\to$}}1
    {←}{{$\leftarrow$}}1
    {×}{{$\times$}}1
    {≥}{{$\geq$}}1
    {≤}{{$\leq$}}1
    {≠}{{$\neq$}}1
    {∈}{{$\in$}}1
    {⌈}{{$\lceil$}}1
    {⌉}{{$\rceil$}}1
    {∞}{{$\infty$}}1
    {α}{{$\alpha$}}1
    {β}{{$\beta$}}1
}
\lstdefinelanguage{TypeScript}{
  keywords={const,let,var,function,return,if,else,for,while,class,interface,type,extends,implements,import,export,from,as,async,await,new,this,true,false,null,undefined,void,any,unknown,never,string,number,boolean,readonly,private,public,protected,enum},
  sensitive=true,
  comment=[l]{//},
  morecomment=[s]{/*}{*/},
  morestring=[b]",
  morestring=[b]',
  morestring=[b]`
}
\lstset{style=codestyle}

\newcommand{\Given}{\textbf{\color{specBlue}Given} }
\newcommand{\When}{\textbf{\color{specBlue}When} }
\newcommand{\Then}{\textbf{\color{specBlue}Then} }
\providecommand{\And}{}
\renewcommand{\And}{\textbf{\color{specBlue}And} }
 (adjust depth).

\documentclass[11pt,a4paper]{article}

\usepackage[utf8]{inputenc}
\usepackage[T1]{fontenc}
\usepackage{lmodern}
\usepackage[a4paper,margin=2.3cm]{geometry}
\usepackage{parskip}
\usepackage{microtype}
\usepackage[dvipsnames]{xcolor}
\usepackage{enumitem}
\usepackage{listings}
\usepackage{tcolorbox}
\usepackage{titlesec}
\usepackage{fancyhdr}
\usepackage{array}
\usepackage{longtable}
\usepackage{booktabs}
\usepackage{amsmath}
\usepackage{amssymb}
\usepackage{hyperref}
\usepackage{cleveref}

\tcbuselibrary{skins,breakable}

\definecolor{specBlue}{RGB}{30,64,132}
\definecolor{specGray}{RGB}{90,90,90}
\definecolor{specAccent}{RGB}{198,40,40}
\definecolor{specGreen}{RGB}{30,110,60}

\hypersetup{
  colorlinks=true,
  linkcolor=specBlue,
  urlcolor=specBlue,
  citecolor=specBlue,
  pdfborder={0 0 0}
}

\titleformat{\section}{\Large\bfseries\color{specBlue}}{\thesection}{0.75em}{}
\titleformat{\subsection}{\large\bfseries\color{specBlue!80}}{\thesubsection}{0.6em}{}

\makeatletter
\newcommand{\specID}[1]{\def\@specID{#1}}
\newcommand{\specTitle}[1]{\def\@specTitle{#1}}
\newcommand{\specStatus}[1]{\def\@specStatus{#1}}
\newcommand{\specVersion}[1]{\def\@specVersion{#1}}
\newcommand{\specOwner}[1]{\def\@specOwner{#1}}
\newcommand{\specTraces}[1]{\def\@specTraces{#1}}

\specID{SPEC-XXX}
\specTitle{Untitled Spec}
\specStatus{DRAFT}
\specVersion{0.1}
\specOwner{Jeremie Nunez}
\specTraces{}

\newcommand{\makespecheader}{%
  \begin{center}
    {\LARGE\bfseries\color{specBlue}\@specTitle}\\[0.4em]
    {\small\color{specGray}\texttt{\@specID} \quad v\@specVersion \quad \textsc{\@specStatus} \quad \@specOwner}
  \end{center}
  \vspace{0.4em}
  \hrule
  \vspace{0.8em}
}

\pagestyle{fancy}
\fancyhf{}
\fancyhead[L]{\small\color{specGray}\@specID}
\fancyhead[R]{\small\color{specGray}\@specTitle}
\fancyfoot[C]{\small\color{specGray}\thepage}
\renewcommand{\headrulewidth}{0pt}
\makeatother

\newtcolorbox{invariant}{
  colback=specBlue!5, colframe=specBlue!75,
  title=Invariant, fonttitle=\bfseries, breakable,
  boxrule=0.5pt, arc=2pt
}

\newtcolorbox{scenario}[1]{
  colback=white, colframe=specBlue!50,
  title={Scenario: #1}, fonttitle=\bfseries, breakable,
  boxrule=0.5pt, arc=2pt
}

\newtcolorbox{ac}[1]{
  colback=specAccent!5, colframe=specAccent!60,
  title={Acceptance Criterion #1}, fonttitle=\bfseries, breakable,
  boxrule=0.5pt, arc=2pt
}

\newtcolorbox{nongoal}{
  colback=specGray!8, colframe=specGray!60,
  title=Non-Goal, fonttitle=\bfseries, breakable,
  boxrule=0.5pt, arc=2pt
}

\lstdefinestyle{codestyle}{
  basicstyle=\ttfamily\small,
  breaklines=true,
  showstringspaces=false,
  frame=single,
  framesep=4pt,
  rulecolor=\color{specBlue!30},
  backgroundcolor=\color{specBlue!2},
  xleftmargin=0.5em,
  xrightmargin=0.5em,
  keywordstyle=\color{specBlue}\bfseries,
  commentstyle=\color{specGray}\itshape,
  stringstyle=\color{specGreen}
}
\lstset{
  inputencoding=utf8,
  extendedchars=true,
  literate=
    {—}{{--}}1
    {–}{{-}}1
    {…}{{...}}1
    {’}{{'}}1
    {‘}{{'}}1
    {“}{{``}}1
    {”}{{''}}1
    {é}{{\'e}}1
    {è}{{\`e}}1
    {ê}{{\^e}}1
    {à}{{\`a}}1
    {â}{{\^a}}1
    {ç}{{\c{c}}}1
    {ô}{{\^o}}1
    {→}{{$\to$}}1
    {←}{{$\leftarrow$}}1
    {×}{{$\times$}}1
    {≥}{{$\geq$}}1
    {≤}{{$\leq$}}1
    {≠}{{$\neq$}}1
    {∈}{{$\in$}}1
    {⌈}{{$\lceil$}}1
    {⌉}{{$\rceil$}}1
    {∞}{{$\infty$}}1
    {α}{{$\alpha$}}1
    {β}{{$\beta$}}1
}
\lstdefinelanguage{TypeScript}{
  keywords={const,let,var,function,return,if,else,for,while,class,interface,type,extends,implements,import,export,from,as,async,await,new,this,true,false,null,undefined,void,any,unknown,never,string,number,boolean,readonly,private,public,protected,enum},
  sensitive=true,
  comment=[l]{//},
  morecomment=[s]{/*}{*/},
  morestring=[b]",
  morestring=[b]',
  morestring=[b]`
}
\lstset{style=codestyle}

\newcommand{\Given}{\textbf{\color{specBlue}Given} }
\newcommand{\When}{\textbf{\color{specBlue}When} }
\newcommand{\Then}{\textbf{\color{specBlue}Then} }
\providecommand{\And}{}
\renewcommand{\And}{\textbf{\color{specBlue}And} }
 (adjust depth).

\documentclass[11pt,a4paper]{article}

\usepackage[utf8]{inputenc}
\usepackage[T1]{fontenc}
\usepackage{lmodern}
\usepackage[a4paper,margin=2.3cm]{geometry}
\usepackage{parskip}
\usepackage{microtype}
\usepackage[dvipsnames]{xcolor}
\usepackage{enumitem}
\usepackage{listings}
\usepackage{tcolorbox}
\usepackage{titlesec}
\usepackage{fancyhdr}
\usepackage{array}
\usepackage{longtable}
\usepackage{booktabs}
\usepackage{amsmath}
\usepackage{amssymb}
\usepackage{hyperref}
\usepackage{cleveref}

\tcbuselibrary{skins,breakable}

\definecolor{specBlue}{RGB}{30,64,132}
\definecolor{specGray}{RGB}{90,90,90}
\definecolor{specAccent}{RGB}{198,40,40}
\definecolor{specGreen}{RGB}{30,110,60}

\hypersetup{
  colorlinks=true,
  linkcolor=specBlue,
  urlcolor=specBlue,
  citecolor=specBlue,
  pdfborder={0 0 0}
}

\titleformat{\section}{\Large\bfseries\color{specBlue}}{\thesection}{0.75em}{}
\titleformat{\subsection}{\large\bfseries\color{specBlue!80}}{\thesubsection}{0.6em}{}

\makeatletter
\newcommand{\specID}[1]{\def\@specID{#1}}
\newcommand{\specTitle}[1]{\def\@specTitle{#1}}
\newcommand{\specStatus}[1]{\def\@specStatus{#1}}
\newcommand{\specVersion}[1]{\def\@specVersion{#1}}
\newcommand{\specOwner}[1]{\def\@specOwner{#1}}
\newcommand{\specTraces}[1]{\def\@specTraces{#1}}

\specID{SPEC-XXX}
\specTitle{Untitled Spec}
\specStatus{DRAFT}
\specVersion{0.1}
\specOwner{Jeremie Nunez}
\specTraces{}

\newcommand{\makespecheader}{%
  \begin{center}
    {\LARGE\bfseries\color{specBlue}\@specTitle}\\[0.4em]
    {\small\color{specGray}\texttt{\@specID} \quad v\@specVersion \quad \textsc{\@specStatus} \quad \@specOwner}
  \end{center}
  \vspace{0.4em}
  \hrule
  \vspace{0.8em}
}

\pagestyle{fancy}
\fancyhf{}
\fancyhead[L]{\small\color{specGray}\@specID}
\fancyhead[R]{\small\color{specGray}\@specTitle}
\fancyfoot[C]{\small\color{specGray}\thepage}
\renewcommand{\headrulewidth}{0pt}
\makeatother

\newtcolorbox{invariant}{
  colback=specBlue!5, colframe=specBlue!75,
  title=Invariant, fonttitle=\bfseries, breakable,
  boxrule=0.5pt, arc=2pt
}

\newtcolorbox{scenario}[1]{
  colback=white, colframe=specBlue!50,
  title={Scenario: #1}, fonttitle=\bfseries, breakable,
  boxrule=0.5pt, arc=2pt
}

\newtcolorbox{ac}[1]{
  colback=specAccent!5, colframe=specAccent!60,
  title={Acceptance Criterion #1}, fonttitle=\bfseries, breakable,
  boxrule=0.5pt, arc=2pt
}

\newtcolorbox{nongoal}{
  colback=specGray!8, colframe=specGray!60,
  title=Non-Goal, fonttitle=\bfseries, breakable,
  boxrule=0.5pt, arc=2pt
}

\lstdefinestyle{codestyle}{
  basicstyle=\ttfamily\small,
  breaklines=true,
  showstringspaces=false,
  frame=single,
  framesep=4pt,
  rulecolor=\color{specBlue!30},
  backgroundcolor=\color{specBlue!2},
  xleftmargin=0.5em,
  xrightmargin=0.5em,
  keywordstyle=\color{specBlue}\bfseries,
  commentstyle=\color{specGray}\itshape,
  stringstyle=\color{specGreen}
}
\lstset{
  inputencoding=utf8,
  extendedchars=true,
  literate=
    {—}{{--}}1
    {–}{{-}}1
    {…}{{...}}1
    {’}{{'}}1
    {‘}{{'}}1
    {“}{{``}}1
    {”}{{''}}1
    {é}{{\'e}}1
    {è}{{\`e}}1
    {ê}{{\^e}}1
    {à}{{\`a}}1
    {â}{{\^a}}1
    {ç}{{\c{c}}}1
    {ô}{{\^o}}1
    {→}{{$\to$}}1
    {←}{{$\leftarrow$}}1
    {×}{{$\times$}}1
    {≥}{{$\geq$}}1
    {≤}{{$\leq$}}1
    {≠}{{$\neq$}}1
    {∈}{{$\in$}}1
    {⌈}{{$\lceil$}}1
    {⌉}{{$\rceil$}}1
    {∞}{{$\infty$}}1
    {α}{{$\alpha$}}1
    {β}{{$\beta$}}1
}
\lstdefinelanguage{TypeScript}{
  keywords={const,let,var,function,return,if,else,for,while,class,interface,type,extends,implements,import,export,from,as,async,await,new,this,true,false,null,undefined,void,any,unknown,never,string,number,boolean,readonly,private,public,protected,enum},
  sensitive=true,
  comment=[l]{//},
  morecomment=[s]{/*}{*/},
  morestring=[b]",
  morestring=[b]',
  morestring=[b]`
}
\lstset{style=codestyle}

\newcommand{\Given}{\textbf{\color{specBlue}Given} }
\newcommand{\When}{\textbf{\color{specBlue}When} }
\newcommand{\Then}{\textbf{\color{specBlue}Then} }
\providecommand{\And}{}
\renewcommand{\And}{\textbf{\color{specBlue}And} }


\specID{SPEC-TH-010}
\specTitle{Nano-Swarm --- parallel cheap retrieval with bounded fan-out}
\specStatus{DRAFT}
\specVersion{0.1}
\specOwner{Jeremie Nunez}

\begin{document}
\makespecheader

\section{Context}
The Explorer subsystem of Thalamus must gather evidence from the open web on every research cycle.
A single call to a mid-tier model with web search returns one perspective and 1--2 domains, which
is too narrow and too monolithic. The Nano-Swarm replaces that single call with a fan-out of many
cheap nano-model workers, each constrained to a distinct \emph{lens} (angle, specialty, or sub-topic).
Fan-out is bounded, executed in fixed-size waves, and every worker is independent, so a partial
failure never blocks the cycle.

\section{Goals}
\begin{itemize}[leftmargin=1.2em]
  \item Turn a small set of high-level exploration queries into a bounded set of micro-queries, each paired with a distinct retrieval lens.
  \item Execute the micro-queries in parallel waves with a fixed concurrency cap and an inter-wave delay to respect upstream rate limits.
  \item Produce a deduplicated set of retrieval artifacts (one per unique source URL, plus synthetic artifacts when a worker returns text without citations).
  \item Emit deterministic swarm statistics (call counts, success ratio, unique domains, wall time, estimated cost) on every cycle.
  \item Fail open: any single worker failure is absorbed; the cycle returns whatever succeeded.
\end{itemize}

\section{Non-Goals}
\begin{nongoal}
This spec does not cover: (a) generation of the upstream exploration queries (owned by the Scout),
(b) relevance or novelty scoring (owned by the Curator, SPEC-TH-012), (c) persistence of artifacts
to the feed or knowledge store, (d) the typed external source adapters (SPEC-TH-011), and
(e) the choice of underlying nano model or its pricing.
\end{nongoal}

\section{Invariants}
\begin{invariant}
Fan-out is bounded: the number of dispatched micro-queries per cycle never exceeds
\texttt{MAX\_MICRO\_QUERIES}, regardless of the number of upstream queries or lenses.
\end{invariant}

\begin{invariant}
Concurrency is bounded: at most \texttt{WAVE\_SIZE} nano calls are in-flight at any moment;
successive waves are separated by at least \texttt{WAVE\_DELAY\_MS} of wall time (except after the last wave).
\end{invariant}

\begin{invariant}
Every single nano call is bounded in time by \texttt{CALL\_TIMEOUT\_MS} and every error, timeout, or
non-2xx response is converted into a typed failure result --- never a thrown exception that escapes the swarm.
\end{invariant}

\begin{invariant}
Artifact URLs are unique across the returned set after normalization (lowercased host, stripped
tracking params, no trailing slash).
\end{invariant}

\begin{invariant}
If no upstream credentials are configured, the swarm returns an empty artifact set and well-formed stats;
it never crashes the caller.
\end{invariant}

\section{Interface}
Public surface of the swarm. The consumer is the Explorer orchestrator.

\begin{lstlisting}[language=TypeScript]
export interface ExplorationQuery {
  query: string;
  type: "web" | "academic" | "market";
  signal: string;
  priority: number;
  maxDepth: number;
}

export interface RetrievalArtifact {
  url: string;              // unique source URL, or "nano://<workerId>" synthetic
  title: string;
  body: string;             // <= 3000 chars, markdown-stripped
  entities: ExtractedEntities;
  dataPoints: string[];
  sourceQuery: string;      // upstream query that produced this artifact
  depth: 0;
}

export interface SwarmStats {
  totalCalls: number;
  successCalls: number;
  failedCalls: number;
  wallTimeMs: number;
  uniqueDomains: number;
  totalUrls: number;
  totalChars: number;
  estimatedCost: number;
}

export class NanoSwarm {
  crawl(upstream: ExplorationQuery[]): Promise<{
    articles: RetrievalArtifact[];
    urlsCrawled: number;
    stats: SwarmStats;
  }>;
}
\end{lstlisting}

\section{Scenarios}

\begin{scenario}{Nominal fan-out}
\Given a list of 6 upstream exploration queries with mixed \texttt{type} values \\
\When \texttt{NanoSwarm.crawl} is invoked \\
\Then the swarm decomposes them into exactly \texttt{MAX\_MICRO\_QUERIES} micro-queries, each tagged with a distinct lens identifier \\
\And the micro-queries are dispatched in waves of \texttt{WAVE\_SIZE} \\
\And the returned \texttt{articles} contain no duplicate normalized URLs \\
\And \texttt{stats.totalCalls} equals the number of micro-queries dispatched
\end{scenario}

\begin{scenario}{Partial failure absorption}
\Given a cycle where 30\% of nano calls return HTTP errors or time out \\
\When the swarm completes \\
\Then the failed calls are counted in \texttt{stats.failedCalls} \\
\And successful calls still contribute artifacts to the output \\
\And no exception is raised to the caller
\end{scenario}

\begin{scenario}{Credential absence}
\Given the upstream API key is not configured in the environment \\
\When the swarm is invoked \\
\Then every nano call returns a typed failure result with an explicit error marker \\
\And \texttt{articles} is empty \\
\And \texttt{stats.successCalls} is 0 and \texttt{stats.totalCalls} equals the number of dispatched micro-queries
\end{scenario}

\begin{scenario}{Rate-limit backpressure}
\Given a set of 50 micro-queries \\
\When the swarm executes them \\
\Then the total wall time is at least \((\lceil 50 / \texttt{WAVE\_SIZE} \rceil - 1) \times \texttt{WAVE\_DELAY\_MS}\) milliseconds \\
\And at no point are more than \texttt{WAVE\_SIZE} requests in flight concurrently
\end{scenario}

\begin{scenario}{Synthetic artifact fallback}
\Given a worker that returns substantive text but cites no URLs \\
\When results are merged \\
\Then exactly one synthetic artifact is produced with \texttt{url = "nano://<workerId>"} \\
\And that artifact carries the worker's \texttt{sourceQuery} and extracted entities
\end{scenario}

\section{Acceptance Criteria}

\begin{ac}{AC-1}
Given an empty upstream query list, \texttt{crawl} returns \texttt{articles: []}, \texttt{urlsCrawled: 0}, and \texttt{stats.totalCalls: 0}.
\end{ac}

\begin{ac}{AC-2}
The number of dispatched nano calls in a single cycle is always \(\le\) \texttt{MAX\_MICRO\_QUERIES}, verified by intercepting the HTTP transport.
\end{ac}

\begin{ac}{AC-3}
Concurrency never exceeds \texttt{WAVE\_SIZE}: a test that instruments the nano caller records at no tick a live-call count greater than \texttt{WAVE\_SIZE}.
\end{ac}

\begin{ac}{AC-4}
After normalization, no two artifacts in a single output share the same URL.
\end{ac}

\begin{ac}{AC-5}
A 4xx or 5xx HTTP response from a single worker increments \texttt{stats.failedCalls} by 1 and does not reduce the contribution of other workers.
\end{ac}

\begin{ac}{AC-6}
A worker whose text exceeds 200 characters and has 0 extracted URLs produces exactly one synthetic artifact with a \texttt{nano://} URL.
\end{ac}

\begin{ac}{AC-7}
\texttt{stats.wallTimeMs} is populated on every return and is monotonically greater than the elapsed time measured outside the call by no more than 50~ms.
\end{ac}

\section{Traceability}

\begin{center}
\begin{tabular}{@{}lll@{}}
\toprule
AC & Test file & Test name \\
\midrule
AC-1 & \texttt{packages/thalamus/tests/nano-swarm.empty.spec.ts} & \texttt{returns empty on empty input} \\
AC-2 & \texttt{packages/thalamus/tests/nano-swarm.fanout.spec.ts} & \texttt{caps fan-out at MAX\_MICRO\_QUERIES} \\
AC-3 & \texttt{packages/thalamus/tests/nano-swarm.concurrency.spec.ts} & \texttt{never exceeds wave size} \\
AC-4 & \texttt{packages/thalamus/tests/nano-swarm.dedup.spec.ts} & \texttt{deduplicates URLs after normalization} \\
AC-5 & \texttt{packages/thalamus/tests/nano-swarm.failure.spec.ts} & \texttt{absorbs per-worker failures} \\
AC-6 & \texttt{packages/thalamus/tests/nano-swarm.synthetic.spec.ts} & \texttt{produces synthetic artifact when no URLs} \\
AC-7 & \texttt{packages/thalamus/tests/nano-swarm.stats.spec.ts} & \texttt{reports accurate wall time} \\
\bottomrule
\end{tabular}
\end{center}

\section{Open Questions}
\begin{itemize}[leftmargin=1.2em]
  \item Should \texttt{WAVE\_SIZE}, \texttt{WAVE\_DELAY\_MS}, and \texttt{MAX\_MICRO\_QUERIES} be tunable per-cycle (config) rather than compile-time constants?
  \item Should the swarm expose per-wave progress as an async iterator for the UI layer, or is the final batch result sufficient?
  \item Is cost estimation in \texttt{stats} authoritative, or should it be computed downstream from a billing feed?
\end{itemize}

\end{document}
